\documentclass[11pt,a4paper]{article}
\usepackage[T1]{fontenc}
\usepackage{lmodern,amsmath,amssymb,amsthm,mathtools,booktabs,microtype}
\usepackage[margin=26mm]{geometry}
\usepackage[colorlinks=true,linkcolor=blue,citecolor=blue,urlcolor=blue]{hyperref}
\newtheorem{theorem}{Theorem}[section]
\newtheorem{proposition}[theorem]{Proposition}
\newtheorem{lemma}[theorem]{Lemma}
\newtheorem{corollary}[theorem]{Corollary}
\theoremstyle{definition}
\newtheorem{definition}[theorem]{Definition}
\newtheorem{remark}[theorem]{Remark}
\newcommand{\Q}{\mathbb Q}\newcommand{\Z}{\mathbb Z}
\newcommand{\OK}{\mathcal O_K}
\setlength{\emergencystretch}{2em}
\title{Locally Dense Diophantine Sets in the Rational Subfield}
\author{}\date{9 November 2026}
\begin{document}\maketitle
\begin{abstract}
We give an explicit polynomial of degree $13$, height $162$, and
four total coordinates whose existential projection over $\Q(i)$ is
a proper subset of $\Q$. This subset is dense simultaneously in every
finite collection of completions of $\Q$, and the polynomial is locally
soluble for every input in every completion of $\Q(i)$. Nevertheless,
the rational input $-1$ has no global witness. The construction uses
the curve $y^2=x^3+9$ and a rank-zero elliptic obstruction.
We prove a general simultaneous density statement for ratios of
elliptic coordinates under rank stability, and give a bounded-denominator
criterion that would suffice for existential definability of a ring of
integers. Local density and the bounded-denominator criterion are
different requirements: the construction in this paper does not define
all of $\Q$ or the Gaussian integers.
\end{abstract}

\section{An explicit polynomial and its scope}

A subset $S$ of a field $L$ is \emph{Diophantine in $L$} if membership
in $S$ can be expressed by the existence of finitely many solutions of
polynomial equations over $L$. The total number of coordinates includes
the distinguished input. The height of an integer polynomial is the
largest absolute value of its coefficients.

Put $K=\Q(i)$ and define
\begin{align}
 F(W,U)&=U^2W-9W^3-1,\label{eq:F}\\
 G(T,W,V)&=V^2TW-9T^3W^3-1,\label{eq:G}\\
 P(T,W,U,V)&=T\bigl(F(W,U)^2-2G(T,W,V)^2\bigr).\label{eq:P}
\end{align}
Let
\begin{equation}\label{eq:D-definition}
 D=\{t\in K:\exists w,u,v\in K,\ P(t,w,u,v)=0\}.
\end{equation}

\begin{theorem}\label{thm:main}
The following statements hold.
\begin{enumerate}
\item $D$ is a proper subset of $\Q$, with $0,1\in D$ and $-1\notin D$.
\item For every finite nonempty set $S$ of places of $\Q$, the diagonal
image of $D$ is dense in $\prod_{v\in S}\Q_v$.
\item For every place $\mathfrak p$ of $K$ and every $t\in K_{\mathfrak p}$,
the equation $P(t,W,U,V)=0$ has a solution over $K_{\mathfrak p}$.
\item The polynomial $P$ has four total coordinates, degree $13$,
height $162$, and $11$ nonzero monomials.
\end{enumerate}
\end{theorem}

The arithmetic needed for this theorem consists of two small rank-zero
calculations, given in Section~\ref{sec:arithmetic} with reproducible
unconditional descent computations. The other parts of the proof are
given in Sections~\ref{sec:global}--\ref{sec:density}.

Theorem~\ref{thm:main} concerns a proper subset of the rational subfield.
Its coordinate count is therefore not a coordinate bound for a
Diophantine definition of $\Z$, $\Q$, or $\Z[i]$ in $\Q(i)$.
In particular, local solubility for every input and simultaneous
local density do not establish the completeness of an existential
definition of the rational subfield.

\section{The elliptic curves and the arithmetic input}\label{sec:arithmetic}

Consider the elliptic curves
\[
 E:y^2=x^3+9,\qquad
 E':y^2=x^3+81,\qquad
 E^-:y^2=x^3-9.
\]
Their origins are denoted by $O$.

\begin{lemma}\label{lem:groups}
One has
\[
 E'(\Q)=\{O,(0,9),(0,-9)\},\qquad E^-(\Q)=\{O\}.
\]
The point $R=(3,6)\in E(\Q)$ has infinite order.
\end{lemma}
\begin{proof}
The rank-zero statements for $E'$ and $E^-$ are the rigorously computed
entries $972b2$ and $3888k1$ in Cremona's tables; the corresponding
entry for $E$ is $972c1$ \cite{CremonaData,LMFDB81,LMFDBReliability}.
They can be reproduced independently by unconditional $2$-descent.
In PARI/GP $2.17.2$, the commands
\begin{verbatim}
ellrank(ellinit([0,0,0,0,81]))
ellrank(ellinit([0,0,0,0,-9]))
\end{verbatim}
both return \verb|[0, 0, 0, []]|. The first two entries are proven
lower and upper bounds on the Mordell--Weil rank. PARI's bounds
use the $2$-Selmer group and the Cassels pairing; no Birch--Swinnerton-Dyer
assumption or numerical analytic-rank calculation is used
\cite[\texttt{ellrank}]{PARI}. The accompanying program
\nolinkurl{check_elliptic_ratio_obstruction.py} repeats these computations.

For completeness, the torsion orders can also be checked by elementary
finite-field calculations. Both curves have good reduction at $5$ and
$7$. Direct counts give
\[
\begin{array}{c|cc}
 &\#E_*(\mathbb F_5)&\#E_*(\mathbb F_7)\\\hline
 E'&6&9\\
 E^-&6&7.
\end{array}
\]
Reduction is injective on torsion of order prime to the residue
characteristic. These two counts bound the torsion order of $E'$ by
$3$ and force the torsion subgroup of $E^-$ to be trivial. The three
displayed points of $E'$ are distinct and rational, proving the asserted
groups once the ranks are zero.

Finally,
\[
                 [2](3,6)=\left(-\frac{15}{16},\frac{183}{64}\right).
\]
The Nagell--Lutz theorem says that a rational torsion point on an
integral short Weierstrass equation has integral coordinates
\cite[Chapter VIII, Section 7]{Silverman}. Thus $(3,6)$ cannot have
finite order.
\end{proof}

\begin{lemma}\label{lem:exact-stability}
Every point of $E(\Q(i))$ is rational:
\[
                          E(\Q(i))=E(\Q).
\]
\end{lemma}
\begin{proof}
Let $\sigma$ be complex conjugation and take $A\in E(\Q(i))$.
Put $B=A-\sigma A$. Then $\sigma B=-B$.
If $B\ne O$, write $B=(a,ib)$ with $a,b\in\Q$: the abscissa is
fixed by $\sigma$ and the ordinate is negated. Its equation gives
\[
              -b^2=a^3+9,\qquad b^2=(-a)^3-9.
\]
Consequently $(-a,b)$ is an affine rational point of $E^-$,
contrary to Lemma~\ref{lem:groups}. Therefore $B=O$,
so $A=\sigma A$ and $A\in E(\Q)$.
\end{proof}

This proof establishes equality of point sets, which is stronger than
equality of ranks. There is no unexamined finite-index quotient in the
definition used below.

\begin{proposition}\label{prop:quadratic-fields}
Let $d\ne1$ be a squarefree integer. If
\[
               E^{(d)}:Y^2=X^3+9d^3
\]
has trivial rational Mordell--Weil group, then
$E(\Q(\sqrt d))=E(\Q)$. The hypothesis holds for
\[
                   d\in\{-15,-13,-11,-10,-1,3,5,7\}.
\]
\end{proposition}
\begin{proof}
With $\sigma$ the nonidentity automorphism, an affine point
$A-\sigma A$ has the form $(a,b\sqrt d)$, with $a,b\in\Q$.
Its equation $db^2=a^3+9$ gives the rational point
$(da,d^2b)$ on $E^{(d)}$. Triviality of $E^{(d)}(\Q)$ thus forces
$A=\sigma A$. For each of the eight displayed values of $d$,
the accompanying checker computes matching unconditional rank
bounds zero and torsion order one for $E^{(d)}$.
\end{proof}

\section{The global projection and the excluded value}\label{sec:global}

Define the set of nonzero abscissas
\[
       X=\{x(A):A\in E(\Q),\ A\ne O,\ x(A)\ne0\}.
\]

\begin{proposition}\label{prop:ratio-identification}
The set in \eqref{eq:D-definition} is exactly
\[
                         D=\{0\}\cup X/X.
\]
\end{proposition}
\begin{proof}
The element $2$ is not a square in $\Q(i)$: if
$(a+bi)^2=2$ with $a,b\in\Q$, then $ab=0$, and neither
$a^2=2$ nor $-b^2=2$ is possible. Hence for $t\ne0$,
\eqref{eq:P} vanishes if and only if $F=G=0$.

The equation $F=0$ forces $w\ne0$ and is equivalent to the assertion
that
\[
                      A=\left(\frac1w,\frac uw\right)
\]
lies on $E$. Likewise $G=0$ forces $tw\ne0$ and is equivalent to
\[
                      B=\left(\frac1{tw},\frac v{tw}\right)\in E.
\]
By Lemma~\ref{lem:exact-stability}, both points are rational.
Their abscissas are nonzero and $x(A)/x(B)=t$, so $t\in X/X$.

Conversely, if $t=x(A)/x(B)$ for points with nonzero abscissas,
put $w=1/x(A)$, $u=y(A)/x(A)$, and $v=y(B)/x(B)$.
These rational coordinates satisfy $F=G=0$.
The factor $T$ in \eqref{eq:P} supplies the input $t=0$.
\end{proof}

The chart used in the proof has a simple geometric interpretation.
In the projective equation $Y^2Z=X^3+9Z^3$, take the affine chart
$X\ne0$ and put $u=Y/X$, $w=Z/X$. The equation becomes $F=0$.
Thus nonvanishing of the original abscissa is enforced by the affine
chart itself, without a separate inverse coordinate. The equation
$G=0$ results from eliminating the second chart coordinate by
$w_2=tw$.

\begin{corollary}\label{cor:same-polynomial}
The same polynomial $P$ defines the same set $D$ in each field
$\Q(\sqrt d)$ listed in Proposition~\ref{prop:quadratic-fields}.
\end{corollary}
\begin{proof}
In each listed field, $2$ is nonsquare: for a squarefree $d\ne1$,
the only quadratic field containing $\sqrt2$ is the one with
$d=2$, which is absent from the list. The proof of
Proposition~\ref{prop:ratio-identification} therefore applies
using Proposition~\ref{prop:quadratic-fields} in place of
Lemma~\ref{lem:exact-stability}.
\end{proof}

\begin{proposition}\label{prop:minus-one}
One has $-1\notin D$.
\end{proposition}
\begin{proof}
Otherwise there would be rational points with abscissas $a$ and $-a$,
where $a\ne0$. If their ordinates are $b$ and $c$, then
\[
          b^2=9+a^3,\qquad c^2=9-a^3.
\]
Multiplying gives
\[
                    (bc)^2=81-a^6=(-a^2)^3+81.
\]
Hence $(-a^2,bc)\in E'(\Q)$. Lemma~\ref{lem:groups} forces its
abscissa to be zero, a contradiction.
\end{proof}

For the size assertion, direct expansion of \eqref{eq:P} gives
\begin{align*}
 P={}&TU^4W^2-18TU^2W^4-2TU^2W+81TW^6+18TW^3-T\\
     &-2T^3V^4W^2+36T^5V^2W^4+4T^2V^2W
       -162T^7W^6-36T^4W^3.
\end{align*}
The eleven monomials are distinct. Their maximum total degree is
$13$, attained by $T^7W^6$, and the maximum coefficient magnitude
is $162$. The supplied sparse expansion checker verifies these
claims with integer arithmetic.

\section{Solubility over every completion}\label{sec:local}

\begin{proposition}\label{prop:all-local}
Let $L$ be $\mathbb R$, $\mathbb C$, or a finite extension of
$\Q_p$. For every $t\in L$ the equation $P(t,W,U,V)=0$ has an
$L$-solution.
\end{proposition}
\begin{proof}
The assertion is immediate for $t=0$. Suppose $t\ne0$.
Squares in $L$ contain a neighborhood of the nonzero square $9$.
For the nonarchimedean fields this follows from Hensel's lemma,
and for the archimedean fields it follows from the local square-root
function at $9$.

Choose $s\in L^\times$ sufficiently close to zero that both
$9+s^3$ and $9+t^3s^3$ are squares. This gives points of $E(L)$
with nonzero abscissas $s$ and $ts$, whose ratio is $t$.
The chart substitution in Proposition~\ref{prop:ratio-identification}
therefore gives a solution with $F=G=0$.
\end{proof}

In particular, the excluded global input $-1$ has witnesses over
every completion. There are uniform elementary choices for its
local abscissa $s$: use $s=p$ for odd $p\ne3$, $s=9$ for $p=3$,
$s=2$ for $p=2$, and $s=1$ over $\mathbb R$.
For $p\ne2,3$ the two radicands reduce to $9$ modulo $p$;
for $p=3$ their quotients by $9$ are $1$ modulo $3$;
for $p=2$ the radicands are $17$ and $1$.
Thus no local obstruction detects this global exclusion.

\section{Simultaneous density of global ratios}\label{sec:density}

Local solubility alone does not prove density of the global values.
The following lemma supplies the required global approximation.

\begin{lemma}\label{lem:cyclic-density}
Let $E/\Q$ be an elliptic curve, let $R\in E(\Q)$ have infinite
order, and let $S$ be a finite set of places of $\Q$.
There is a positive integer $N$ such that the closure of
$\{[Nn]R:n\in\Z\}$ in $\prod_{v\in S}E(\Q_v)$ contains a product
of open neighborhoods of $O$.
\end{lemma}
\begin{proof}
For each finite $p\in S$, choose an open subgroup of $E(\Q_p)$
on which the elliptic logarithm is an analytic group isomorphism
onto $p^{a_p}\Z_p$. After choosing $N$ sufficiently divisible,
$[N]R$ belongs to all these subgroups. Its logarithm $\ell_p$ is
nonzero, because the logarithm is injective there and $R$ is
not torsion. The closure of its integer multiples is therefore
the open subgroup corresponding to $\ell_p\Z_p$.

The Chinese remainder theorem shows that the diagonal integers
are dense in $\prod_{p\in S,\ p<\infty}\Z_p$. Thus the integer
multiplier can approximate independently any prescribed element
of each $\ell_p\Z_p$.

If the real place belongs to $S$, enlarge $N$ so that $[N]R$ lies
in $E(\mathbb R)^0$. Identify this connected real Lie group with
$\mathbb R/\mathbb Z$. The image $\alpha$ of $[N]R$ is irrational,
since otherwise the point would be torsion. Any finite list of
$p$-adic approximation requirements on $n$ defines an arithmetic
progression $n=n_0+Mk$. Along this progression,
$n\alpha$ is dense modulo one because $M\alpha$ is irrational.
The real approximation can therefore be imposed simultaneously
with the finite-place conditions. This proves the assertion.
\end{proof}

\begin{proposition}\label{prop:general-density}
Let $E/\Q$ have positive rank, and let $C\in E(\Q)$ be an affine
point that is not a point of order two. Put $c=x(C)$ and
\[
 X_c=\{x(A)-c:A\in E(\Q),\ A\ne O,\ x(A)\ne c\}.
\]
Then $X_c/X_c$ is dense in $\prod_{v\in S}\Q_v$ for every finite
nonempty set $S$ of places of $\Q$.
\end{proposition}
\begin{proof}
Fix a point $R\in E(\Q)$ of infinite order and apply
Lemma~\ref{lem:cyclic-density}. The rational points
$C+[Nn]R$ are dense in a product of neighborhoods of $C$.
The function $f=x-c$ has nonzero derivative at $C$, since
$dx=2y\,\omega$ for the invariant differential $\omega$, and
$y(C)\ne0$. The local inverse function theorem therefore shows
that the values $f(C+[Nn]R)$ are dense in a product
$\prod_{v\in S}U_v$, where each $U_v$ is a neighborhood of zero
in $\Q_v$.

Take target values $a_v\in\Q_v$. For each $v$, choose
$b_v\in U_v\setminus\{0\}$ sufficiently small that $a_vb_v\in U_v$.
Choose independently two global points in the above sequence,
with their $f$-values sufficiently close to $(a_vb_v)_v$ and
$(b_v)_v$. Both values can be required to be nonzero: the zeros
of $f$ are only $C$ and $-C$, so at most finitely many integer
multipliers are excluded. Taking their quotient approximates
$(a_v)_v$ as closely as desired.
\end{proof}

Apply Proposition~\ref{prop:general-density} to the curve $E$ above
and $C=(0,3)$. Then $c=0$ and $X_c=X$. Together with
Propositions~\ref{prop:ratio-identification}, \ref{prop:minus-one},
and \ref{prop:all-local}, this completes the proof of
Theorem~\ref{thm:main}.

\begin{remark}
The density is in completions of $\Q$. The rational subset $D$ is
not asserted to be dense in a completion $K_{\mathfrak p}$ in which
the closure of $\Q$ is a proper local subfield.
Proposition~\ref{prop:all-local} is a different assertion: it allows
the witnesses themselves to range over the completion.
\end{remark}

\section{The general rank-stability construction}

Ratios of local parameters on a stable abelian variety are used for
weak approximation in \cite[Theorem 4.2(b)]{Poonen}. In the elliptic
case the following statement includes the real place and records
an explicit coordinate bound.

\begin{proposition}\label{prop:six-coordinates}
Let $K/\Q$ be a number field and let
$E:y^2=x^3+ax+b$ be an elliptic curve over $\Q$ with
\[
              0<\operatorname{rank}E(\Q)=\operatorname{rank}E(K).
\]
There is a Diophantine subset $D_E\subseteq\Q$ inside $K$ that
is simultaneously dense in every finite product of completions of
$\Q$. It has a definition with six total coordinates and four
equations. If $m\ge1$ annihilates $E(K)/E(\Q)$, the equations have
degree at most $\max(3,3m^2-1)$.
\end{proposition}
\begin{proof}
The quotient $E(K)/E(\Q)$ is finite by the Mordell--Weil theorem
and equality of ranks. Choose $m$ annihilating it and a point
$R\in E(\Q)$ of infinite order. Put $c=x([m]R)$.
Write the abscissa of multiplication by $m$ as
\[
                 x([m]A)=\frac{\phi_m(x(A))}{B_m(x(A))},
\]
where $\deg\phi_m=m^2$ and $\deg B_m=m^2-1$; these are the
usual division-polynomial expressions, with $B_1=1$.
Put $A_m=\phi_m-cB_m$.

For coordinates $(t,x_1,y_1,x_2,y_2,z)$ impose
\begin{align*}
 y_j^2-x_j^3-ax_j-b&=0\qquad(j=1,2),\\
 tA_m(x_2)B_m(x_1)-A_m(x_1)B_m(x_2)&=0,\\
 zB_m(x_1)B_m(x_2)A_m(x_2)-1&=0.
\end{align*}
The last equation excludes every denominator and the zero
denominator of the ratio. Each multiplied point is rational, so
the resulting input $t$ is rational. The degree bounds are
$3$, $2m^2$, and $3m^2-1$, respectively.

To prove simultaneous density, use rational inputs near $R$ in
the sequences provided by Lemma~\ref{lem:cyclic-density}.
The function $x\circ[m]-c$ has a simple zero at $R$, because
the isogeny $[m]$ is unramified in characteristic zero and
$[m]R$ is not two-torsion. The proof of
Proposition~\ref{prop:general-density} applies unchanged.
\end{proof}

This proposition supplies locally rich rational values. It does not
assert that every rational number is a value. Theorem~\ref{thm:main}
shows that completeness can fail even when rank stability is exact
and every local fiber is soluble.

\section{Fibers and rational parametrizations}

\begin{proposition}\label{prop:genus-four}
For $t\in\Q\setminus\{0,1\}$, the smooth projective curve associated
to
\[
                 y^2=s^3+9,\qquad z^2=t^3s^3+9
\]
has genus four. In particular, the nonzero-ratio witnesses for each
such fixed rational input form a finite set.
\end{proposition}
\begin{proof}
Over an algebraic closure, the function field is a biquadratic
extension of $\overline\Q(s)$. Each square root has three finite
branch points and is ramified at infinity. The two finite branch
sets are disjoint because $t^3\ne1$. Thus there are seven branch
points in their composite, each with ramification index two.
Riemann--Hurwitz gives
\[
                     2g-2=4(-2)+7\cdot2=6,
\]
so $g=4$. Faltings's theorem gives finiteness of its rational points
\cite{Faltings}. The affine nonzero-ratio witnesses are a subset.
\end{proof}

The points with $s=0$ supply four rational points on the displayed
projective curve. They are precisely the points removed when
forming a nonzero abscissa ratio. For $t=-1$, the obstruction in
Proposition~\ref{prop:minus-one} excludes every remaining affine
point with $s\ne0$.

\begin{proposition}\label{prop:no-section}
Let $k$ be a characteristic-zero field, let $E/k$ be an elliptic
curve, and let $h\in k(E^r)$ for some $r\ge1$. The rational map
$h:E^r\dashrightarrow\mathbb P^1$ has no rational section.
In particular, there are no $w,u,v\in k(T)$ satisfying $F=G=0$
as identities for the polynomials in \eqref{eq:F}--\eqref{eq:G}.
\end{proposition}
\begin{proof}
Every point of $E(k(T))$ is constant: a rational map
$\mathbb P^1_k\dashrightarrow E$ extends to a morphism by
properness, and a nonconstant morphism is impossible by
Riemann--Hurwitz. Indeed, after extending scalars to an algebraic
closure its ramification formula would read $-2\ge0$.

A putative rational section of a ratio therefore has constant
elliptic points and hence constant ratio. For the explicit system,
$F=G=0$ would give the two points
\[
             (1/w,u/w),\qquad(1/(Tw),v/(Tw))
\]
over $k(T)$, both constant, while their abscissa ratio is the
nonconstant parameter $T$. This is impossible.
\end{proof}

This obstruction concerns rational parametrizations. It does not
exclude an arithmetic procedure that chooses unrelated rational
points separately for different inputs.

\section{A sufficient bounded-denominator condition for integrality}

The following criterion gives a different target for the number-field
problem. Unlike local density, it directly suffices for an existential
definition of the full ring of integers.

\begin{theorem}[Integral-basis criterion; cf.\ {\cite[Theorem 4.2(d)]{Poonen}}]\label{thm:sandwich}
Let $K/\Q$ have degree $n$, and suppose that $S\subset K$ is
Diophantine in $K$ with
\[
                           \Z\subseteq S\subseteq\OK.
\]
Then $\OK$ is Diophantine in $K$.
If $S$ has a definition with $N$ total coordinates, one polynomial,
and degree $d$, then $\OK$ has a definition with $nN$ total
coordinates and $n$ equations of degree at most $d$.
\end{theorem}
\begin{proof}
Choose an integral basis $\omega_1,\ldots,\omega_n$ of $\OK$.
The two inclusions in the hypothesis give
\[
                     \OK=\omega_1S+\cdots+\omega_nS.
\]
Indeed, the right side is integral, and the integer coefficients
of every integral-basis expansion belong to $S$.

For an input $x$, retain $s_1,\ldots,s_{n-1}$ and eliminate
$s_n=(x-\sum_{j<n}\omega_js_j)/\omega_n$.
Impose $n$ copies of the definition of $S$, with disjoint witness
blocks. The count is $1+(n-1)+n(N-1)=nN$; the substitutions
are linear and do not increase degree.
\end{proof}

\begin{corollary}\label{cor:bounded-denominators}
The conclusion still holds if $\Z\subseteq S$ and
$MS\subseteq\OK$ for some positive rational integer $M$.
With the same notation, $nN$ total coordinates and $n$ equations
of degree at most $Md$ suffice.
\end{corollary}
\begin{proof}
The finite union
\[
                 S'=\bigcup_{j=0}^{M-1}(j+MS)
\]
lies in $\OK$ and contains $\Z$, by division with remainder.
If $f(T,\mathbf Y)=0$ defines $S$, the polynomial
\[
           \prod_{j=0}^{M-1}M^d
                  f\!\left(\frac{T-j}{M},\mathbf Y\right)
\]
defines $S'$ using the same witness block. A product over a field
vanishes if and only if one factor vanishes, so shared witnesses
create no additional condition. Its degree is at most $Md$.
Apply Theorem~\ref{thm:sandwich}.
\end{proof}

Neither criterion requires that $S$ be a ring. Neither follows from
the simultaneous density proved above. In fact, the set $D$ in
Theorem~\ref{thm:main} has unbounded denominators at every prime,
as its density in each $\Q_p$ immediately shows.

\section{Counterexample searches and the distinction between ambient rings}

The explicit exclusion in this paper is not a counterexample to
the existence of an existential definition of $\Z[i]$ in $\Q(i)$.
It excludes one proposed way of obtaining the complete rational
subfield from stable elliptic coordinates.

The available undefinability results for large, or ample, fields
do not provide such a number-field counterexample. A field is
large if every smooth curve with a rational point has infinitely
many rational points. Fehm's theorem and Kwon's extension show
that infinite Diophantine sets in a perfect large field cannot
be confined to a proper subfield or even to a finite union of
affine translates of proper subfields \cite{Fehm,Kwon}.
No number field is large: a smooth genus-two curve such as
$y^2=x^5-x+1$, with point $(0,1)$, has only finitely many points
over each number field by Faltings's theorem. The polynomial is
squarefree, since a common root of $x^5-x+1$ and $5x^4-1$
would have to be $5/4$, which does not satisfy the derivative
equation.

Similarly, the rank-stability transfer theorem of
Mazur--Rubin--Shlapentokh defines one ring of integers inside
another ring of integers; its witnesses are integral
\cite[Theorem 1.1 and Section 14]{MRS}. Its approximate-logarithm
method does not by itself remove that witness-domain restriction.
For an approach through elliptic ratios, the remaining requirement
is a proven finite construction that contains every rational
number, or a construction satisfying the bounded-denominator
criterion above. Local approximation alone supplies neither.

\section{Reproducible computation}

The program \nolinkurl{check_elliptic_ratio_obstruction.py} requires
\texttt{cypari2} and uses PARI's unconditional rank bounds. It also
checks the torsion orders, counts the reductions at $5$ and $7$
by two independent finite-field methods, verifies the displayed
doubling formula, expands \eqref{eq:P} with exact integer
coefficients, and records the local square conditions.
The result is stored in
\nolinkurl{elliptic-ratio-certificate.json}.
The commands
\begin{verbatim}
python3 -m pip install cypari2
python3 check_elliptic_ratio_obstruction.py \
  --output elliptic-ratio-certificate.json
\end{verbatim}
reproduce the calculation. The rank-zero bounds, rather than a
finite search for points, are the decisive computational inputs.

\begin{thebibliography}{99}
\bibitem{CremonaData}
J. E. Cremona, \emph{Elliptic curve data}, entries $972b2$, $972c1$,
and $3888k1$ in \texttt{allcurves.00000-09999},
\url{https://github.com/JohnCremona/ecdata}.

\bibitem{LMFDB81}
The LMFDB Collaboration, \emph{Elliptic curve $972.d2$ (Cremona
label $972b2$)},
\url{https://www.lmfdb.org/EllipticCurve/Q/972/d/2}.

\bibitem{LMFDBReliability}
The LMFDB Collaboration, \emph{Reliability of elliptic curve data
over $\Q$},
\url{https://www.lmfdb.org/EllipticCurve/Q/Reliability}.

\bibitem{PARI}
The PARI Group, \emph{PARI/GP reference documentation: Elliptic curves},
especially \texttt{ellrank} and \texttt{elltors},
\url{https://pari.math.u-bordeaux.fr/dochtml/ref/Elliptic_curves.html}.

\bibitem{Silverman}
J. H. Silverman, \emph{The Arithmetic of Elliptic Curves},
second edition, Graduate Texts in Mathematics 106,
Springer, 2009.

\bibitem{Faltings}
G. Faltings, Endlichkeitss\"atze f\"ur abelsche Variet\"aten
\"uber Zahlk\"orpern, \emph{Invent. Math.} \textbf{73} (1983),
349--366.

\bibitem{Fehm}
A. Fehm, Subfields of ample fields. Rational maps and definability,
\emph{J. Algebra} \textbf{323} (2010), 1738--1744.

\bibitem{Kwon}
A. Kwon, \emph{A note on Diophantine subsets of large fields},
arXiv:2411.03212, 2024,
\url{https://arxiv.org/abs/2411.03212}.

\bibitem{MRS}
B. Mazur, K. Rubin, and A. Shlapentokh,
Existential definability and diophantine stability,
\emph{J. Number Theory} \textbf{254} (2024), 1--64;
corrigendum, \textbf{262} (2024), 539--540;
arXiv:2208.09963v3,
\url{https://arxiv.org/abs/2208.09963}.

\bibitem{Poonen}
B. Poonen, \emph{Rank stability makes rings of integers diophantine},
3 July 2026,
\url{https://math.mit.edu/~poonen/papers/h10_over_OK.pdf}.
\end{thebibliography}
\end{document}
